\chapter{Specifikáció}
\label{chap:elemzes}

Ebben a fejezetben a diplomaterv kitűzött céljait bontom le konkrét műszaki követelményekre. A feladat nem csupán egy működő szimuláció elkészítése, hanem egy olyan keretrendszer kialakítása, amely szoftvertechnológiai szempontból is megállja a helyét.

\section{Funkcionális követelmények}

A rendszernek a következő alapvető funkciókat kell ellátnia:

\begin{itemize}
    \item \textbf{Ágensalapú szimulációs környezet:} A Unity3D motoron belül egy lehatárolt, 3D-s élettér biztosítása, ahol az egyedek interakcióba léphetnek egymással és a környezettel.
    \item \textbf{Többszintű populációdinamika:} Legalább két faj (préda és ragadozó) jelenléte, amelyek életciklusuk során táplálkoznak, szaporodnak és elpusztulnak.
    \item \textbf{Genetikai öröklődés modellje:} Az egyedek tulajdonságainak (sebesség, látótávolság, energiafelhasználás) genotípus alapú tárolása és mutációval kombinált átörökítése az utódokba.
    \item \textbf{Eseménykezelés és naplózás:} A szimuláció során bekövetkező releváns események (születés, halál, vadászat) detektálása és strukturált formában történő rögzítése.
    \item \textbf{Külső adatkapcsolat:} A kinyert adatok HTTP protokollon keresztül történő továbbítása egy külső tároló/elemző egység felé.
\end{itemize}

Fontos kiemelni, hogy a fenti követelményekből a BSc projekt során az ágensalapú környezet alapjai és egy kezdeti préda-populáció (nyulak) modellje készült el. Ebben a fázisban már megjelentek az örökölhető tulajdonságok, azonban a rendszer architektúrája még monolitikus volt, és a szimuláció kizárólag egyetlen faj túlélési stratégiáit vizsgálta zárt környezetben.

A jelenlegi munka során az alábbi követelmények teljesítése jelenti az előrelépést:
\vspace{0.3cm}
\begin{itemize}
    \item A \textbf{ragadozó faj (rókák)} bevezetése, amely megteremti a valódi populációdinamikai interakciókat.
    \item A rendszer \textbf{teljes refaktorálása} moduláris, eseményvezérelt architektúrává, megszüntetve a korábbi kódminta strukturális hiányosságait.
    \item \textbf{Logolás és a külső adatbázis-kapcsolat (InfluxDB, Grafana)} implementálása, amely korábban statikus fájl alapú vagy hiányzó funkció volt.
\end{itemize}

\section{Nem-funkcionális követelmények}

A szoftver minőségére vonatkozó elvárások:

\begin{itemize}
    \item \textbf{Kiterjeszthetőség (Extensibility):} Az architektúrának lehetővé kell tennie új fajok vagy viselkedésminták hozzáadását a meglévő kód bázis minimális módosításával (Open/Closed Principle).
    \item \textbf{Teljesítmény, skálázhatóság és valós idejűség:} A szimulációnak stabil képkockasebesség mellett kell futnia több száz aktív ágens jelenléte esetén is. Ehhez szükséges az állapotellenőrzések (polling) kiváltása eseményvezérelt logikával, biztosítva, hogy a statisztikai adatok gyűjtése és külső adatbázisba továbbítása ne okozzon észrevehető akadást vagy késleltetést a vizuális megjelenítésben.
    \item \textbf{Karbantarthatóság:} A BSc szakaszban azonosított „god class” struktúra felszámolása, a felelősségi körök szétválasztása (Single Responsibility Principle).
\end{itemize}

\section{A szimulációs paraméterek specifikációja}

A pontosabb elemzés érdekében rögzíteni kell azokat a változókat, amelyeket a rendszernek kezelnie kell. A szimuláció bemeneti adatait és a vizsgált változók rendszerezett listáját a \ref{table:parameters}. táblázat foglalja össze.

\begin{table}[h!]
\centering
\begin{tabular}{|l|l|}
\hline
\textbf{Kategória} & \textbf{Paraméterek} \\ \hline
Környezet & Terület mérete, táplálék (fű) újratermelődési rátája. \\ \hline
Egyed (Genetika) & Mozgási sebesség, látómező szöge/távolsága, éhségküszöb. \\ \hline
Populáció & Kezdeti egyedszám, mutációs ráta. \\ \hline
Rendszer & Logolási gyakoriság, szimulációs időlépték (Timescale). \\ \hline
\end{tabular}
\caption{A szimuláció kulcsparaméterei}
\label{table:parameters}
\end{table}

A táblázatban felsorolt paraméterek finomhangolása alapvetően meghatározza a szimulált ökoszisztéma stabilitását. Például a mutációs ráta és a kezdeti populációméret közötti egyensúly kritikus a természetes szelekció modellezésekor.

\section{Felhasználói esetek (Use Cases)}

A rendszer elsődleges felhasználója a kutató/hallgató, aki a következő műveleteket végzi:
\begin{enumerate}
	\item  \textbf{Konfiguráció:} A szimuláció indítása előtt beállítja a kezdőparamétereket.
	\item \textbf{Megfigyelés:} Valós időben követi a vizuális térben zajló folyamatokat.
	\item \textbf{Elemzés:} A futás után vagy alatt a generált grafikonok és statisztikák alapján következtetéseket von le a populáció életképességéről.
\end{enumerate}
% SF: ha ki tudod egészíteni kisebb al use case-ekkel a fentieket, akkor egy use case diagram is készíthető a fentiekből, egyébként mivel kevés use case van, nem szükséges